\documentclass[12pt]{article}
\usepackage[utf8]{inputenc}
\usepackage[T1]{fontenc}
\usepackage{hyperref}
\title{Sake Language Resources for K-12 Curriculum Development}
\author{Compiled by Claude Sonnet 5.5 (claude-sonnet-5-5) for Andr\'as Kornai}
\date{}
\begin{document}
\maketitle

\section{Introduction}
Sake (ISO 639-3 \texttt{sak}; Glottocode sake1247) is a Bantu language of Gabon. Glottolog-based sources list it in the Atlantic-Congo family, under Ngomic and the Sake-Ndambomo group. The speaker count is about 8,000 (a 2007 figure). Wikipedia describes the language as undocumented and threatened. Joshua Project lists the alternate names Asake and Shake. It gives a population of about 17,000 for the Sake people, who live mainly in Ogoou\'e-Ivindo Province around Boou\'e. The searches found no school materials or documentation of the language.

\section{Primary Educational Resources}
No primers, textbooks or school programmes for Sake turned up in the searches.

\section{Linguistic and Reference Materials}
\begin{itemize}
\item Wikipedia entry on the Sake language: \url{https://en.wikipedia.org/wiki/Sake_language}
\item Glottolog languoid page (classification, sake1247): \url{https://glottolog.org/resource/languoid/id/sake1247.bigmap.html}
\item ASJP language page (a word list): \url{https://asjp.clld.org/languages/SAKE}
\item Verbix endangered-languages page, which notes the Ethnologue and Guthrie classification differences: \url{https://docs.verbix.com/EndangeredLanguages/Sake}
\end{itemize}

\section{Religious and Community Resources}
Joshua Project has pages on the Sake people and the language, which give demographic information only: \url{https://joshuaproject.net/people_groups/14652} and \url{https://joshuaproject.net/languages/sak}. No scripture translations or community organisations were found.

\section{Digital and Audio Resources}
No audio, video or corpus resources were found.

\section{Curriculum Development Context}
Sake is essentially undocumented, and no teaching materials were found. Any curriculum would first need new documentation work: word lists, a grammar sketch, an orthography and recordings made with older speakers. French is the school language of Gabon.

\end{document}